\documentclass[12pt]{article}
%\usepackage{amssymb}
%\usepackage{amsthm}
%\usepackage{amsgen}
\usepackage{amsfonts}

\newtheorem{theorem}{Theorem} 
\newtheorem{corollary}[theorem]{Corollary} 
\newtheorem{conjecture}[theorem]{Conjecture} 
\newtheorem{lemma}[theorem]{Lemma} 
\newtheorem{proposition}[theorem]{Proposition} 
\newtheorem{definition}[theorem]{Definition} 
\newtheorem{example}[theorem]{Example} 
\newtheorem{exercise}[theorem]{Exercise} 
\newtheorem{axiom}{Axiom} 
\newtheorem{remark}{Remark} 
\def\ppp{{\mathbb{P}}}
\def\aaa{{\mathbb{A}}}
\def\fff{{\mathbb{F}}}
\def\qqq{\mathbb{Q}}
\def\rrr{\mathbb{R}}
\def\ccc{\mathbb{C}}
\def\zzz{\mathbb{Z}}
\def\nnn{\mathbb{N}}
\def\pf{{\bf proof}:\ }
\def\qed{$\Box$}

\begin{document}

\author{David Joyner}
\title{Representations of finite groups on the Riemann-Roch
space of a toric variety}
\date{6-23-2002}
\maketitle

\begin{abstract}
We study the action 
of a finite group on the Riemann-Roch
space of a toric variety. Our main result is the following:
if $G$ is a finite subgroup of the (toric)
automorphism group of a complete, smooth toric variety $X$ and if $D$ is a 
divisor on $X$ left stable by $G$ then we can write the natural representation
of $G$ on $L(D)$ is a direct sum of explicitly described 
irreducible subrepresentations. We also determine explicitly the 
trace of this representation.
\end{abstract}

\tableofcontents

\vskip .2in

Let $X$ be a smooth projective variety over an algebraically
closed field $F$ and let $G$ be a finite subgroup of automorphisms of $X$
over $F$. We often identify $X$ with its set of
$F$-rational points $X(F)$.
If $D$ is a divisor
\footnote{Since $X$ is smooth, we need not distinguish between 
Weil divisors and Cartier divisors (Hartshorne \cite{H},
Prop. II.6.11, and Remark II.6.11.1A).} 
of $X$ which $G$ leaves 
fixed then $G$ acts on the Riemann-Roch space $L(D)$.
We ask: which (modular) representations 
arise in this way? If $\pi$ is such a representation,
what is its trace?

In the case 
when $G$ acts on a curve $X$, this question has been
investigated by many mathematicians - Hurwitz,
Chevalley and Weil, among others. (We refer, for example, to
\cite{JT} for further references.)
In this short note, we discuss the simpler case of 
a group acting (as toric automorphisms) on a toric variety.


\section{The action of $G$ on $L(D)$}

Let $X$ be a smooth projective variety over an algebraically closed
field $F$. Let $F(X)$ denote the function field of
$X$ (the field of rational functions on $X$) and, if
$D$ is any divisor on $X$, let
\[
L(D)= \{f\in F(X)^\times \ |\ (f)+D\geq 0\}\cup \{0\},
\]
and let $\ell(D)$ denote its dimension.
This definition makes sense even if $X$ is 
merely normal and $D$ is a Weil divisor.

The action of ${\rm Aut}(X)$ on $F(X)$ is 
defined by 
\[
\begin{array}{cccc}
\rho:&{\rm Aut}(X)&\rightarrow &{\rm Aut}(F(X)),\\
 & g &\longmapsto & (f\longmapsto f^g)
\end{array}
\]
where $f^g(x)=(\rho(g)(f))(x)=f(g^{-1}(x))$.

Of course, ${\rm Aut}(X)$ also acts on the group $Div(X)$ of
divisors of $X$, denoted $g:D\longmapsto g(D)=\sum_P d_Pg(P)$,
for $g\in {\rm Aut}(X)$ and $D=\sum_P d_P P\in Div(X)$.
It is easy to show that $(f^g)=g(f)$, where
$(f)$ denotes the (principal) divisor of the function
$f\in F(X)$. Because of this, if $(f)+D\geq 0$
then $(f^g)+g(D)\geq 0$, for all $g\in {\rm Aut}(X)$.

If $G\subset {\rm Aut}(X)$ leaves $D\in Div(X)$ fixed then we denote
the induced representation of $G$ in $L(D)$ again
by $\rho$:
\[
\rho:G\rightarrow {\rm Aut}(L(D)).
\]

\section{An example}
\label{sec:example}

Let $\Delta$ be a fan in a lattice $L$. Denote by
$\tau_1$, ..., $\tau_n$ the edges or rays of
the fan and let $v_i$ denote the first (smallest)
lattice point along the ray $\tau_i$.
Let $D_i$ denote the Weil divisor 
\[
D_i={\rm Hom}(\tau_i^\perp\cap L^\perp,\ccc^\times),
\]
which maybe regarded as the closure of the orbit of $T$ acting on 
the edge $\tau_i$. 

Let $\Delta$ be the fan generated by 
\[
v_1=2e_1-e_2,\ \ v_2=-e_1+2e_2,\ \ v_3=-e_1-e_2.
\]

In the notation above, the divisor $D=d_1D_1+d_2D_2+d_3D_3$
is a Cartier divisor
\footnote{This is an Exercise on page 65 of \cite{F},
the solution of which is an easy calculation using the
Exercise on page 62, which is in turn, basically
solved in the back of the book.}
if and only if
$d_1\equiv d_2\equiv d_3\ ({\rm mod}\ 3)$.

\begin{picture}(256.00,150.00)(-120.00,0.00)
\thicklines
\put(0.00,64.00){\circle*{4}} %dot
\put(0.00,128.00){\circle*{4}} 

\put(-64.00,0.00){\circle*{4}} 
\put(0.00,-64.00){\circle*{4}} 

\put(-64.00,-64.00){\circle*{4}} 
\put(-64.00,64.00){\circle*{4}} 
\put(-64.00,128.0){\circle*{4}} 

\put(64.00,0.00){\circle*{4}} 
\put(128.00,0.00){\circle*{4}} 
\put(192.00,0.00){\circle*{4}} 
%\thinlines

\put(0.00,0.00){\vector(-1,2){64.00}} 
\put(0.00,0.00){\vector(-1,-1){64.00}} 

\put(64.00,64.00){\circle*{4}} 
\put(128.00,0.00){\circle*{4}} 
\put(0.00,128.00){\circle*{4}} 
\put(-128.00,0.00){\circle*{4}} 

\put(128.00,80.00){$\sigma_1$} 
\put(-128.00,64.00){$\sigma_2$} 
\put(50.00,-100.00){$\sigma_3$} 

\put(135.00,-75.00){$v_1$} 
\put(-75.00,135.00){$v_2$} 
\put(-75.00,-75.00){$v_3$} 

\put(0.00,-64.00){\circle*{4}} 
\put(64.00,-64.00){\circle*{4}} 
\put(128.00,-64.00){\circle*{4}} 

\put(0.00,0.00){\vector(2,-1){129.00}} %basis vector
\end{picture}

\vskip 2in

Let $G$ denote the automorphism group of $\Delta$ 
generated by the map $g$ sending $(x,y)$ to $(y,x)$, swapping 
$v_1$ and $v_2$ and leaving $v_3$ fixed. By Theorem
1.13 in Oda \cite{O}, 
this corresponds to a $T$-equivariant automorphism of
$X(\Delta)$.

Let 
\[
\begin{array}{c}
P_D=\{(x,y)\ |\ \langle (x,y),v_i\rangle \geq -d_i, \ \forall i\}\\
=\{(x,y)\ |\ 2x-y \geq -d_1, -x+2y \geq -d_2, -x-y \geq -d_3 
\}
\end{array}
\]
denote the polytope associated to the Weil
divisor $D=d_1D_1+d_2D_2+d_3D_3$,
where $D_i$ is as above. 

Let $d_1=d_2=6$ and $d_3=0$.
Then $P_D$ is a triangle in the plane 
with vertices at $(-6,-6)$, $(-2,2)$, and $(2,-2)$. Note that it remains
invariant under the action of $G$. Moreover, the $G$ action
on $P_D\cap L^\perp$ has $7$ singleton orbits
(the lattice points $(-i,-i)$, where $0\leq i\leq 6$)
and $12$ orbits of size $2$.

In this example, the patch $U_{\sigma_1}$ is an affine variety with 
coordinates $x_1,x_2,x_3$ given by $x^3-x_2x_3=0$. 
The automorphism $g$ acts on $U_{\sigma_1}$ sending
$(x_1,x_2,x_3)$ to $(x_1,x_3,x_2)$. The torus embedding
$T\hookrightarrow U_{\sigma_1}$ is given by
sending $(t_1,t_2)$ to 
$(x_1,x_2,x_3)=(t_1t_2,t_1t_2^2,t_1^2t_2)$.

The patch $U_{\sigma_2}$ is an affine variety with 
coordinates $y_1,y_2,y_3$ given by $y^2-y_1y_3=0$. 
The automorphism $g$ does not act on $U_{\sigma_2}$.
The torus embedding
$T\hookrightarrow U_{\sigma_2}$ is given by
sending $(t_1,t_2)$ to 
$(y_1,y_2,y_3)=(t_1^{-2}t_2^{-1},t_1^{-1},t_1^{-1}t_2)$.

The patch $U_{\sigma_3}$ is an affine variety with 
coordinates $z_1,z_2,z_3$ given by $z^2-z_1z_3=0$. 
The automorphism $g$ does not act on $U_{\sigma_3}$.
The torus embedding
$T\hookrightarrow U_{\sigma_3}$ is given by
sending $(t_1,t_2)$ to 
$(x_1,x_2,x_3)=(t_1^{-1}t_2^{-2},t_2^{-1},t_1t_2^{-1})$.
The automorphism $g$ sends $U_{\sigma_2}$ to $U_{\sigma_3}$
by sending $(y_1,y_2,y_3)$ to $(z_1,z_2,z_3)$.


\section{The main result}

The following is our main result.
Let $\Delta$ be a fan associated to an integral
lattice $L$ in $\rrr^n$. Let $P_D$ denote the polytope
associated to the Weil divisor $D$ on $X=X(\Delta)$.


\begin{theorem}
Let $X=X(\Delta)$ be a complete toric variety defined over
an algebraically closed field $F$. Suppose $G\subset {\rm Aut}(X)$
is a finite subgroup, either ${\rm char}(F)=0$ or
$p={\rm char}(F)$ does not divide $|G|$, and that the 
Weil divisor $D$ on $X$ is fixed by $G$.
Let $P_D$ denote the polytope
associated to the Weil divisor $D$ on $X$ and let
$S=P_D\cap L^\perp$. Let 
\[
S=S_1\cup S_2\cup ...\cup S_k
\]
denote the decomposition of $S$ into disjoint primitive $G$-orbits.
For each $m\geq 1$, the natural representation $\pi$ of $G$ on
$L(D)$ decomposes into a direct sum of $k$ irreducible
permutation representations
\[
\pi =\oplus_{i=1}^k \pi_i,
\]
where $\pi_i$ is $|S_i|$-dimensional.
\end{theorem}

\pf
First, note that there is a natural isomorphism
$L(D)\cong \Gamma(X,{\cal O}(D))$
(see for example Griffiths and Harris \cite{GH}, page 136).
We use this map to pull back the action of $G$ on
$\Gamma(X,{\cal O}(D))$ to an action on $L(D)$.
By Fulton \cite{F}, page 66, we have
\[
L(D)=\oplus_{u\in S} F\cdot \chi^u,
\]
where $\chi=(x_1,...,x_n)$ and $\chi^u$ is the
associated monomial in multi-index notation.
This shows that $\pi$,
with representation space $V=L(D)$ 
has an $|S_i|$-dimensional subrepresentation $\pi_i$,
with representation space $V_i=\oplus_{u\in S_i} F\cdot \chi^u$.
If $F$ has characteristic $p$ and $p$ does not divide $|G|$
then every finite dimensional representation of $G$
is semi-simple (Mascke's Theorem, Thrm 3.14, \cite{CR}, or
\cite{Se}, \S 15.7).
From these, the theorem follows.
\qed

Question: Given $G$, $X$, and $D$ as above. Do the
(equivalence classes of) representations of $G$
occurring in $L(D)$ determine $D$ up to linear equivalence?

Let $(\dots ,\dots )$ denote the inner product
on $V=L(D)=\oplus_{u\in S} F\cdot \chi^u$
given by
\[
(\sum_{u\in S}a_u\chi^u,\sum_{v\in S}b_v\chi^v)=
\sum_{w\in S}a_wb_w.
\]
What is the trace of the representation $\pi$ in 
the above theorem?

\begin{theorem}
Let $X$, $G$, $D$, $S$ and $F$ be as in the above theorem.
Assume, in addition, $F$ has characteristic $0$.
Then 
\[
{\rm tr}(\pi(g))=|\{u\in S\ |\ g(u)=u\}|=|Fix_S(g)|.
\]
\end{theorem}

\pf
This follows from well-known facts on permutation representations.
\qed

For example, the trace of the representation $\pi$ in 
\S \ref{sec:example} is 
\[
{\rm tr}(\pi(g))=
\left\{
\begin{array}{cc}
31, & g=1\\
7, & g\not= 1.
\end{array}
\right.
\]



\begin{thebibliography}{99}  

\bibitem[CR]{CR} C. Curtis, I. Reiner, {\bf Methods of 
representation theory, I}, Wiley-Interscience, 1981.
 
\bibitem[F]{F} W. Fulton, {\bf Introduction to toric varieties},
Annals of Math. Studies, Princeton Univ. Press, 1993.

\bibitem[GH]{GH} P. Griffiths and J. Harris, 
{\bf Principles of algebraic geometry}, John Wiley and Sons,
1978.

\bibitem[H]{H} R. Hartshorne, {\bf Algebraic geometry},
Springer-Verlag, 1977.

\bibitem[JT]{JT} D. Joyner and W. Traves, 
``Representations of finite groups on the Riemann-Roch
space of a curve,'', available at
\newline
\verb++

\bibitem[O]{O} T. Oda, {\bf Convex bodies and algebraic geometry},
Springer-Verlag, 1985.

\bibitem[Se]{Se} J.-P. Serre, {\bf Linear representations of 
finite groups}, Springer-Verlag, 1977.

%\bibitem[VJ]{VJ} H. Verrill and D. Joyner, ``Notes on toric varieties,''
%available at
%\newline
%\verb++

\end{thebibliography}

\end{document}


